% Provenance (never printed): NSI premiere, "Representation des donnees: types et valeurs de base" (writing an
% integer in base b; binary representation of a relative integer, two's complement; size in bits of a sum or a
% product). Bitwise operations and the 2038 problem are additions. See sources/book2/BRIEFS.md.
\chapter{Integers in the Machine}\label{ch:g10:machine-integers}

Noor's grandfather kept an old step counter clipped to his belt. One summer he walked across a whole region, and
each evening he wrote down the number it showed. On the eighth day the number was smaller than on the seventh,
although he had walked all day. The counter was not broken. It had six digits on its little screen and, inside, a
memory of a fixed number of bits; after the largest number it could hold, it went back to zero, like a car's
odometer.

Every number a processor works on lives in a fixed number of bits. This chapter shows how whole numbers, positive
and negative, are written in $8$, $16$, $32$ or $64$ bits, what happens when a result does not fit, and how a
program can work on the bits of a number one by one.

\begin{recall}
A whole number is written in base $2$ with the digits $0$ and $1$, each position worth twice the one on its right;
in base $16$ with the digits $0$ to $9$ and $A$ to $F$; Python's \code{bin}, \code{hex} and \code{int(s, 2)}
convert between bases, and two binary numbers are added column by column with a carry (Book 1, grade 8). A bit
is one binary digit and a byte is eight bits (Book 1, grade 5). $2^{10} = 1024$, and $2^{a} \times 2^{b} =
2^{a+b}$ (mathematics book, grade 8).
\end{recall}

\section{Unsigned integers on \texorpdfstring{$n$}{n} bits}\label{sec:g10:machine-integers:unsigned}

Base $2$ is one base among many. In any base $b \ge 2$, a whole number is written with digits from $0$ to $b - 1$,
and the digit at position $i$ (counting from $0$ on the right) is worth $b^i$ times its value. The digits are found
from the right by repeated division: the remainder of $n$ divided by $b$ is the last digit, and the quotient holds
the others.

\omcode{code/grade-10/03-machine-integers/listings/to_base.py}{1}{17}{Writing a number in any base from $2$ to
$36$.}\label{lst:g10:machine-integers:to-base}

\omoutput{code/grade-10/03-machine-integers/outputs/to_base.txt}{The same number in five bases. \code{int} reads a
string of digits back in a given base.}\label{out:g10:machine-integers:to-base}

A processor does not write numbers with as many digits as needed: its registers and memory cells have a fixed
size.

\begin{definition}[Unsigned integer, \texorpdfstring{$n$}{n}-bit integer]\label{def:g10:machine-integers:unsigned}
An \emph{$n$-bit integer}\index{n-bit integer@$n$-bit integer} is a whole number stored in exactly $n$ bits,
$b_{n-1} \dots b_1 b_0$. Read as an \emph{unsigned integer}\index{unsigned integer}, its value is
$b_{n-1} 2^{n-1} + \dots + b_1 2 + b_0$. The usual sizes are $8$ bits (one byte), $16$, $32$ and $64$ bits.
\end{definition}

\begin{proposition}[Range of \texorpdfstring{$n$}{n} bits]\label{prop:g10:machine-integers:range}
The \omterm{def:g10:machine-integers:unsigned}{unsigned integers} on $n$ bits are exactly the whole numbers from $0$ to $2^n - 1$. A whole number $x \ge 1$ needs
$k$ bits, where $k$ is the number with $2^{k-1} \le x < 2^k$.
\end{proposition}

\begin{proof}
There are $2$ choices for each of the $n$ bits, so $2^n$ patterns, and different patterns have different values
(base $2$ writing is unique). The smallest value is $0$ (all bits $0$). The largest has all bits $1$: adding $1$ to
$11\dots1$ carries through every column and gives $100\dots0$, which is $2^n$; so $11\dots1$ is $2^n - 1$. Hence the
$2^n$ values are $0, 1, \dots, 2^n - 1$. A number with $2^{k-1} \le x < 2^k$ fits in $k$ bits but not in $k - 1$,
whose largest value is $2^{k-1} - 1$.
\end{proof}

So one byte holds $0$ to $255$, $16$ bits hold $0$ to $65{,}535$, and $32$ bits $0$ to $4{,}294{,}967{,}295$.
% ledger: g10:py-bit-length
Python's method \code{n.bit_length()} gives the number of bits that $n$ needs.

\omcode{code/grade-10/03-machine-integers/listings/bits_needed.py}{1}{7}{How many bits a number, a sum and a product
need.}\label{lst:g10:machine-integers:bits-needed}

\omoutput{code/grade-10/03-machine-integers/outputs/bits_needed.txt}{Each power of two needs one more bit than the
number just below it.}\label{out:g10:machine-integers:bits-needed}

\begin{proposition}[Bits of a sum and of a product]\label{prop:g10:machine-integers:sum-product}
If $a$ fits in $p$ bits and $b$ in $q$ bits, with $p \ge q$, then $a + b$ fits in $p + 1$ bits and $a \times b$ fits in
$p + q$ bits.
\end{proposition}

\begin{proof}
$a \le 2^p - 1$ and $b \le 2^q - 1 \le 2^p - 1$, so $a + b \le 2^{p+1} - 2 < 2^{p+1}$. And
$a \times b \le (2^p - 1)(2^q - 1) < 2^p \times 2^q = 2^{p+q}$.
\end{proof}

The run above shows that both bounds can be reached: the largest numbers on $20$ and $13$ bits have a sum of $21$
bits and a product of $33$.

When a result needs more bits than the machine has, the extra bits are simply lost.

\begin{definition}[Overflow]\label{def:g10:machine-integers:overflow}
An \emph{overflow}\index{overflow} happens when the result of an operation does not fit in the number of bits that
hold it. An $n$-bit unsigned adder then keeps the $n$ lowest bits of the true result, that is, the remainder of the
true result divided by $2^n$: the value \emph{wraps around}.
\end{definition}

Python's \texttt{\%} gives that remainder, so a few lines simulate the step counter, here with $16$ bits.

\omcode{code/grade-10/03-machine-integers/listings/wrap.py}{1}{13}{A 16-bit counter, and one addition on $8$ and on
$16$ bits.}\label{lst:g10:machine-integers:wrap}

\omoutput{code/grade-10/03-machine-integers/outputs/wrap.txt}{After $65{,}535$ comes $0$: a whole walk's worth of
steps has vanished.}\label{out:g10:machine-integers:wrap}

The toy machine of the next year (grade 11) has registers of one byte. The same addition, $250 + 10$, run on it:

\omcode{code/grade-10/03-machine-integers/listings/toy_wrap.toy}{1}{3}{Three instructions for the toy
machine.}\label{lst:g10:machine-integers:toy-wrap}

\omoutput{code/grade-10/03-machine-integers/outputs/toy_wrap.txt}{Register \code{R0} holds $4$, and the flag
\code{C} (carry) is $1$: the machine kept the lost bit in a flag.}\label{out:g10:machine-integers:toy-wrap}

The processor does not stop at an \omterm{def:g10:machine-integers:overflow}{overflow}; it records it in its carry flag, and a program that does not look at
the flag never knows.

\section{Signed integers: two's complement}\label{sec:g10:machine-integers:signed}

Negative numbers need a sign, but a memory cell holds only bits. The usual method gives the leftmost bit a negative
weight.

\begin{definition}[Two's complement]\label{def:g10:machine-integers:twos-complement}
In \emph{two's complement}\index{two's complement} on $n$ bits, the pattern $b_{n-1} b_{n-2} \dots b_0$ stands for
the signed value
\[
  -b_{n-1} 2^{n-1} + b_{n-2} 2^{n-2} + \dots + b_1 2 + b_0 .
\]
The leftmost bit is the \emph{sign bit}: it is $1$ exactly for the negative values.
\end{definition}

\begin{omfigure}
\begin{tikzpicture}[font=\small]
  \foreach \w [count=\i] in {$-128$,$64$,$32$,$16$,$8$,$4$,$2$,$1$} {
    \node[cell small, minimum width=11mm] (b\i) at (1.1*\i,0) {1};
    \node[font=\footnotesize, above=1mm of b\i] {\w};
  }
  \fill[omProp!25] (b1.south west) rectangle (b1.north east);
  \node[cell small, minimum width=11mm] at (b1) {1};
  \node[anchor=west, font=\footnotesize, align=left] at (9.6,0)
    {$-128 + 64 + 32 + 16 + 8 + 4 + 2 + 1$\\$= -1$};
  \node[anchor=east, font=\footnotesize] at (0.45,0.48) {weights};
\end{tikzpicture}

{\small The weights of the bits of one byte in \omterm{def:g10:machine-integers:twos-complement}{two's complement}. Only the leftmost weight is negative; the byte
of eight $1$ bits stands for $-1$.}
\end{omfigure}

The first half of the patterns, those starting with $0$, keep their unsigned value; the second half stands for
negative numbers. The listing prints the sixteen patterns of $4$ bits.

\omcode{code/grade-10/03-machine-integers/listings/twos.py}{1}{18}{Encoding and decoding $n$-bit \omterm{def:g10:machine-integers:twos-complement}{two's
complement}.}\label{lst:g10:machine-integers:twos}

\omoutput{code/grade-10/03-machine-integers/outputs/twos.txt}{On $4$ bits, from $-8$ to $7$. The unsigned and the
signed reading differ by $16$ in the second half.}\label{out:g10:machine-integers:twos}

\begin{proposition}[Range and negation]\label{prop:g10:machine-integers:negation}
The signed integers on $n$ bits are the whole numbers from $-2^{n-1}$ to $2^{n-1} - 1$. The pattern of $x$ and its
unsigned reading $u$ satisfy $u = x$ if $x \ge 0$ and $u = x + 2^n$ if $x < 0$. To negate a number, flip every bit and
add $1$ (keeping $n$ bits).
\end{proposition}

\begin{proof}
If the sign bit is $0$, the value is that of the other $n - 1$ bits, from $0$ to $2^{n-1} - 1$. If it is $1$, the
value is $-2^{n-1}$ plus a number from $0$ to $2^{n-1} - 1$, so from $-2^{n-1}$ to $-1$; and it is the unsigned
value, which counts the leftmost bit as $+2^{n-1}$, minus $2 \times 2^{n-1} = 2^n$. For the negation, write $\bar x$
for the pattern with every bit of $x$ flipped. In each column exactly one of $x$ and $\bar x$ has a $1$, so the
unsigned readings add up to $11\dots1 = 2^n - 1$. Hence $\bar x + 1 = 2^n - x$, which on $n$ bits is the pattern of
$-x$.
\end{proof}

\begin{example}[\texorpdfstring{$-5$}{-5} on 8 bits]\label{ex:g10:machine-integers:minus-five}
$5$ is \code{00000101}; flipping every bit gives \code{11111010}; adding $1$ gives \code{11111011}. Check:
$-128 + 64 + 32 + 16 + 8 + 2 + 1 = -5$.
\end{example}

\begin{omfigure}
\begin{tikzpicture}[font=\footnotesize]
  \draw[thick, black!40] (0,0) circle (2.2);
  \foreach \u in {0,...,15} {
    \pgfmathsetmacro{\ang}{90 - 22.5*\u}
    \pgfmathtruncatemacro{\s}{\u < 8 ? \u : \u - 16}
    \node[circle, fill=white, draw, inner sep=1pt, minimum size=5mm] at (\ang:2.2) {\s};
    \node[font=\scriptsize\ttfamily, text=black!60] at (\ang:3.0) {\u};
  }
  \draw[very thick, omProp, -{Stealth}] (90-22.5*7+8:1.55) arc (90-22.5*7+8:90-22.5*8-8:1.55);
  \node[omProp, align=center] at (0,-0.9) {overflow:\\$7 + 1$ gives $-8$};
  \draw[very thick, omDef, -{Stealth}] (90-22.5*15+8:1.55) arc (90-22.5*15+8:90-22.5*16-8:1.55);
  \node[omDef, align=center] at (0,1.0) {no overflow:\\$-1 + 1$ gives $0$};
  \node[text=black!60, align=left, anchor=west] at (3.3,0.6) {grey: the unsigned\\reading of the pattern};
\end{tikzpicture}

{\small The sixteen patterns of $4$ bits on a wheel, with their signed values. Adding $1$ moves one step clockwise.
The signed \omterm{def:g10:machine-integers:overflow}{overflow} is the step from $7$ to $-8$; the unsigned one is the step from $15$ (that is, $-1$) to $0$.}
\end{omfigure}

The great advantage of \omterm{def:g10:machine-integers:twos-complement}{two's complement} is that the addition does not change. Since a negative $x$ is stored as
$x + 2^n$, and the adder works modulo $2^n$, the unsigned adder of the previous section adds signed numbers
correctly as long as the true result is in the range. One circuit serves both; grade 11 builds it from gates.

When the true signed result is out of range, the result wraps around the wheel and the sign comes out wrong:
$100 + 50$ on $8$ bits gives the pattern of $150$, which reads as $-106$. This can be spotted without the true
result.

\begin{proposition}[Detecting a signed overflow]\label{prop:g10:machine-integers:signed-overflow}
The sum of two signed $n$-bit integers \omterm{def:g10:machine-integers:overflow}{overflows} exactly when the two operands have the same sign and the result's
sign bit is different.
\end{proposition}

The proof is \cref{exo:g10:machine-integers:14}.

\section{Working on the bits}\label{sec:g10:machine-integers:bitwise}

\begin{definition}[Bitwise operation]\label{def:g10:machine-integers:bitwise}
A \emph{bitwise operation}\index{bitwise operation} works on each bit position separately. In Python,
\code{a & b} has a $1$ where both \code{a} and \code{b} have a $1$ (\textsc{and}); \code{a | b} where at least one
has a $1$ (\textsc{or}); \code{a ^ b} where exactly one has a $1$ (exclusive or, \textsc{xor}); and \code{~a} flips
every bit (\textsc{not}; for Python's integers, \code{~a} equals \code{-a - 1}).
\end{definition}

\begin{definition}[Shift]\label{def:g10:machine-integers:shift}
A \emph{shift}\index{shift (bit shift)} moves every bit of a number by $k$ positions: \code{a << k} to the left,
adding $k$ zeros on the right, and \code{a >> k} to the right, dropping the $k$ rightmost bits. For a whole number
$a \ge 0$, \code{a << k} is $a \times 2^k$ and \code{a >> k} is the quotient \code{a // 2**k}.
\end{definition}

The most common use of these operations is to pick out, switch on or switch off one bit, using a second number that
has $1$ bits only at the places of interest.

\begin{definition}[Mask]\label{def:g10:machine-integers:mask}
A \emph{mask}\index{mask (bit mask)} is a number used with a \omterm{def:g10:machine-integers:bitwise}{bitwise operation} to select some bits of another number:
\code{x & m} keeps the bits of \code{x} where \code{m} has a $1$ and clears the others; \code{x | m} sets them to
$1$; \code{x ^ m} flips them. The mask \code{1 << i} has a single $1$, at position $i$.
\end{definition}

\begin{omfigure}
\begin{tikzpicture}[font=\small\ttfamily]
  \def\row#1#2#3#4{\node[anchor=east, font=\footnotesize\rmfamily] at (-0.2,#1) {#2};
    \foreach \c [count=\i] in {#3} \node[cell small, #4] at (0.6*\i,#1) {\c};}
  \row{0}{value \code{x}}{1,0,1,0,0,1,1,0}{}
  \row{-0.7}{mask \code{~(1 << 5)}}{1,1,0,1,1,1,1,1}{fill=omDef!10}
  \row{-1.4}{\code{x & mask}}{1,0,0,0,0,1,1,0}{fill=omProp!12}
  \draw[omProp, thick] (0.6*3-0.3,0.35) rectangle (0.6*3+0.3,-1.75);
  \node[anchor=west, font=\footnotesize\rmfamily, align=left] at (5.3,-0.7) {bit 5 cleared,\\the others unchanged};
\end{tikzpicture}

{\small Clearing bit $5$ of \code{x} (bits are numbered from $0$ on the right). The \omterm{def:g10:machine-integers:mask}{mask} has a $0$ at position $5$ and
$1$ elsewhere; \textsc{and} copies every bit of \code{x} except that one.}
\end{omfigure}

\omcode{code/grade-10/03-machine-integers/listings/masks.py}{1}{13}{\omterm{def:g10:machine-integers:mask}{Masks} on a byte of flags, and the three bytes
of a colour.}\label{lst:g10:machine-integers:masks}

\omoutput{code/grade-10/03-machine-integers/outputs/masks.txt}{\code{format(x, "08b")} writes \code{x} in binary on
eight digits. The colour splits into three bytes, and the three bytes give the colour
back.}\label{out:g10:machine-integers:masks}

A colour on a screen is often stored as $24$ bits, one byte for each of red, green and blue (Book 1 showed the
\texttt{\#F4A261} notation of web pages). A \omterm{def:g10:machine-integers:shift}{shift} moves the wanted byte to the right end, and the \omterm{def:g10:machine-integers:mask}{mask} \code{0xFF},
eight $1$ bits, keeps that byte alone.

\begin{omfigure}
\medskip
\begin{bytefield}[bitwidth=1.45em]{24}
  \bitheader[endianness=big]{0,7,8,15,16,23} \\
  \bitbox{8}{red: \code{0xF4} = 244} & \bitbox{8}{green: \code{0xA2} = 162} & \bitbox{8}{blue: \code{0x61} = 97}
\end{bytefield}

{\small The colour \code{0xF4A261} as three bytes. Red is \code{(c >> 16) & 0xFF}, green \code{(c >> 8) & 0xFF},
blue \code{c & 0xFF}.}
\end{omfigure}

\begin{method}[Testing, setting and clearing bit \texorpdfstring{$i$}{i}]\label{met:g10:machine-integers:bit-ops}
\begin{enumerate}
  \item Test: \code{(x >> i) & 1} is $1$ when bit $i$ of \code{x} is $1$.
  \item Set to $1$: \code{x | (1 << i)}.
  \item Clear to $0$: \code{x & ~(1 << i)}.
  \item Flip: \code{x ^ (1 << i)}.
\end{enumerate}
\end{method}

\section{Integers that never overflow}\label{sec:g10:machine-integers:python-int}

% ledger: g10:py-int-unlimited
Python hides all of this: its integers have unlimited precision. When a result does not fit in a machine word,
Python uses as many words as needed, at the price of slower arithmetic on very large numbers.

\omcode{code/grade-10/03-machine-integers/listings/big_int.py}{1}{6}{Large integers in
Python.}\label{lst:g10:machine-integers:big-int}

\omoutput{code/grade-10/03-machine-integers/outputs/big_int.txt}{$2^{100}$ needs $101$ bits; the last line goes past
the largest signed $64$-bit integer without any trouble.}\label{out:g10:machine-integers:big-int}

\begin{inpractice}[Fixed widths in other languages]
% ledger: g10:c-int32
Most other languages, among them C, C++, Rust and Java, use the processor's fixed widths. The C type
\code{int32_t}, for instance, is a signed integer of exactly $32$ bits; adding $1$ to the largest one is an error
that the program may not notice. When code in
those languages computes with large values (money in cents, times in microseconds, file sizes), the first question
is always: how many bits does the largest value need? The university books meet these types.
\end{inpractice}

\begin{history}[The year 2038]
% ledger: g10:posix-epoch, g10:posix-86400
Many systems count time as a number of seconds since an instant called the \emph{Epoch}, fixed by the POSIX
standard at 1~January 1970, 00:00:00 UTC, every day counting exactly $86{,}400$ seconds. A count of seconds kept in a
signed $32$-bit integer runs out in 2038, and the weekend problem computes the exact second. Software written to
count in $64$ bits does not have the problem; old devices and old file formats that store $32$ bits do, which is why
the date is known as the year 2038 problem.
\end{history}

\section{Exercises}\label{sec:g10:machine-integers:exercises}

\begin{exercise}[$\star$]\label{exo:g10:machine-integers:1}
What do \code{bin(13)}, \code{hex(255)}, \code{int("1010", 2)}, \code{int("ff", 16)} and \code{bin(2**8)} return?
\end{exercise}

\begin{exercise}[$\star$]\label{exo:g10:machine-integers:2}
A sensor sends its readings as $12$-bit \omterm{def:g10:machine-integers:unsigned}{unsigned integers}. What are the smallest and the largest values? How many
different values can it send?
\end{exercise}

\begin{exercise}[$\star$]\label{exo:g10:machine-integers:3}
Write $-5$, $-1$, $-128$ and $127$ in \omterm{def:g10:machine-integers:twos-complement}{two's complement} on $8$ bits. Which of $128$ and $-129$ can be written, and why?
\end{exercise}

\begin{exercise}[$\star$]\label{exo:g10:machine-integers:4}
What do \code{1 << 10}, \code{100 >> 2}, \code{7 >> 3}, \code{5 << 1} and \code{0b1011 >> 1} give? Say each time which
multiplication or division by a power of two it computes.
\end{exercise}

\begin{exercise}[$\star$]\label{exo:g10:machine-integers:5}
Without looking at the chapter's listing, write \code{to_base(n, b)} with its \omterm{def:g10:specifications-and-tests:docstring}{docstring} and test it by converting
back with \code{int}. What does it return for $n = 0$, and why must that case be treated apart?
\end{exercise}

\begin{exercise}[$\star\star$]\label{exo:g10:machine-integers:6}
Without computing the product, give the largest number of bits that the product of a $20$-bit number and a $13$-bit
number can need, and the smallest. Give an example for each.
\end{exercise}

\begin{exercise}[$\star\star$]\label{exo:g10:machine-integers:7}
A program reads 8-bit two's-complement values from a file as strings of eight characters \code{0} and \code{1}, and
decodes them with \code{int(bits, 2)}. A temperature of $-5$ degrees appears as $251$. Explain the bug and write a
correct \code{decode(bits)}.
\end{exercise}

\begin{exercise}[$\star\star$]\label{exo:g10:machine-integers:8}
Write \code{test_bit(x, i)}, \code{set_bit(x, i)} and \code{clear_bit(x, i)}, and test that clearing a bit after
setting it gives the same result as clearing it directly.
\end{exercise}

\begin{exercise}[$\star\star$]\label{exo:g10:machine-integers:9}
On $8$ bits, compute $100 + 50$: the pattern of the result, its unsigned value and its signed value. Which of the two
readings \omterm{def:g10:machine-integers:overflow}{overflows}? Do the same for $-100 + (-50)$.
\end{exercise}

\begin{exercise}[$\star\star$]\label{exo:g10:machine-integers:10}
Follow the values of \code{a} and \code{b} through \code{a = a ^ b}, \code{b = a ^ b}, \code{a = a ^ b}, starting with
\code{a = 12} and \code{b = 10}, in binary. What do the three lines do? Use the facts \code{x ^ x == 0} and
\code{x ^ 0 == x} to explain why.
\end{exercise}

\begin{exercise}[$\star\star$]\label{exo:g10:machine-integers:11}
To read the green byte of a colour, a student writes \code{(c >> 16) & 0xFF}; another writes \code{c >> 8}. For
\code{c = 0xF4A261}, what does each get? Correct both.
\end{exercise}

\begin{exercise}[$\star\star\star$]\label{exo:g10:machine-integers:12}
Prove that for every signed $n$-bit integer $x$ other than $-2^{n-1}$, flipping the bits of $x$ and adding $1$ gives the
pattern of $-x$, and say what happens for $-2^{n-1}$ (try it on $4$ bits).
\end{exercise}

\begin{exercise}[$\star\star\star$]\label{exo:g10:machine-integers:13}
The \emph{population count} of $x \ge 0$ is the number of $1$ bits of $x$. Write a function that looks at the bits one
by one with \code{x & 1} and \code{x >> 1}, and a second one that repeats \code{x = x & (x - 1)} until $x$ is $0$.
Explain why \code{x & (x - 1)} clears the rightmost $1$ bit, and compare the numbers of loop turns of the two functions
for $x = 2^{40}$.
\end{exercise}

\begin{exercise}[$\star\star\star$]\label{exo:g10:machine-integers:14}
Prove the proposition on detecting a signed \omterm{def:g10:machine-integers:overflow}{overflow}: if $x$ and $y$ have different signs, $x + y$ is always in range;
if both are $\ge 0$ and the true sum is too large, the stored result has its sign bit set; if both are negative and the
true sum is too small, it has its sign bit clear.
\end{exercise}

\begin{exercise}[$\star\star\star$]\label{exo:g10:machine-integers:15}
Write \code{add_n_bits(a, b, n)}, which adds two $n$-bit patterns and returns the $n$-bit result, the carry (the
unsigned sum needed $n + 1$ bits) and the signed \omterm{def:g10:machine-integers:overflow}{overflow} flag. Test it against Python's own arithmetic for every
pair of $8$-bit values you can afford, and give the result for \code{(100, 50, 8)}.
\end{exercise}

\section{Problem: The 2038 Rollover}\label{sec:g10:machine-integers:problem}

\begin{problem}[{Weekend problem --- the last second of a 32-bit clock}]\label{pb:g10:machine-integers:1}
A door-entry system records the time of each badge as the number of seconds since 1~January 1970, 00:00:00 UTC (see
the history box of this chapter), in a signed $32$-bit integer. The program \code{exercises/rollover.py} and
Python's \code{datetime} module help with the dates, but every count below can be done by hand first.

\needspace{4\baselineskip}\medskip\noindent\textbf{Part I --- Counting seconds.}
\begin{enumerate}
  \item How many seconds are there in a day, and in a year of $365$ days?
  \item How many days are there from 1~January 1970 to 1~January 2000? (There are $7$ leap years among
    1970--1999; name them.)
  \item Deduce the number of seconds since the Epoch at 1~January 2000, 00:00:00 UTC.
  \item Check it with Python: \code{datetime(2000, 1, 1, tzinfo=UTC).timestamp()}, after
    \code{from datetime import UTC, datetime}.
\end{enumerate}

\needspace{4\baselineskip}\medskip\noindent\textbf{Part II --- The largest value.}
\begin{enumerate}[resume]
  \item What is the largest signed $32$-bit integer? Give it as a power of two minus one and in decimal.
  \item Write its pattern in hexadecimal.
  \item About how many years of $365.25$ days does it represent?
  \item Add that to 1970: in which year does the counter run out?
\end{enumerate}

\needspace{4\baselineskip}\medskip\noindent\textbf{Part III --- The exact instant.}
\begin{enumerate}[resume]
  \item Run \code{exercises/rollover.py}. At which date and time does the largest value fall?
  \item How many whole days after the Epoch is that, and how many seconds remain after those days?
  \item Convert the remaining seconds into hours, minutes and seconds, and check the time printed by the program.
  \item What is the next second, written as a count, and why can the register not hold it?
\end{enumerate}

\needspace{4\baselineskip}\medskip\noindent\textbf{Part IV --- After the wrap.}
\begin{enumerate}[resume]
  \item Adding $1$ to the largest value in a signed $32$-bit register gives which pattern, and which signed value?
  \item Using the wheel of this chapter, explain the jump on a $4$-bit clock that counts from $0$.
  \item Which date does the program print for the wrapped value?
  \item By how many seconds does the clock jump backwards? Express it as a power of two.
  \item How long before 1970 is the wrapped date, in days? Check with the answer to question 10.
  \item If the count were stored as an \emph{unsigned} $32$-bit integer instead, when would it run out? (The program
    prints it.) What would the system then be unable to record?
  \item With $64$ signed bits, how long would the count last? (The program prints it.)
  \item State the result: the last instant a signed $32$-bit count of seconds can hold, and the date it shows one
    second later.
\end{enumerate}
\end{problem}
